\documentclass[11pt]{article}
\usepackage[margin=1in]{geometry}
\usepackage[T1]{fontenc}
\usepackage{lmodern,amsmath,amssymb,amsthm,mathtools,hyperref}
\setlength{\emergencystretch}{3em}
\newtheorem{theorem}{Theorem}[section]
\newtheorem{lemma}[theorem]{Lemma}
\newtheorem{corollary}[theorem]{Corollary}
\newtheorem{proposition}[theorem]{Proposition}
\newcommand{\E}{\mathbb E}
\newcommand{\tr}{\operatorname{tr}}
\title{Quantitative partial Hadamard matrices,\\near-square instability, and rational dilations}
\author{Joshua Mark\thanks{ORCID: \href{https://orcid.org/0009-0003-5905-2382}{0009-0003-5905-2382}. AI tools assisted with derivation, independent review, computation and exposition.}}
\date{20 November 2026 --- research draft; priority not established}
\begin{document}\maketitle

\begin{abstract}
The multiplier theorem combines with a precise conditional
prime-distribution estimate to give partial Hadamard matrices with
at least $n/2-\sqrt n\exp(O(\log n/\log\log n))$ rows.
Only Dirichlet GRH for power-of-two conductors is needed for that
result, and it remains explicitly conditional.
Separately, the same unconditional multiplier theorem gives exact
near-optimal isotropic laws with only $n+n^{1/2+o(1)}$ atoms but
distance $1-(1+o(1))n^{-1/2}$ from every Hadamard mixture.
Third completion moments improve two finite certificates.
An exact rational-dilation parameterization and a structured
order-twenty example show why arithmetic concatenations cannot
automatically be completed while preserving their rows.
These results neither construct unknown Hadamard orders nor prove
the full conjecture; bibliographic novelty is unestablished.
\end{abstract}
The unconditional existence inputs are read from the existing project.
No attempt to validate the recent GRH announcement is part of this
pass. Conditional and unconditional statements below are kept separate.
\section{A conditional square-root-scale deficit for partial Hadamard matrices}
\label{sec:conditional-half-frame}

Let $r(n)$ denote the largest number of orthogonal sign rows of length
$n$, and let $\tau(u)$ be the least exponent $t$ for which a Hadamard
matrix of order $2^tu$ exists, for odd $u$. The established multibase
input is unconditional:
\[
 \tau(u)=O\!\left(\frac{\log u}{\log\log u}\right).
\]
We combine it with a \emph{conditional} prime-distribution estimate.
Throughout, logarithms are natural. No assertion that GRH has been
proved or verified is made.

The concatenation strategy is historical: de Launey--Gordon,
\emph{A Comment on the Hadamard Conjecture}, JCTA 95 (2001),
180--184, Lemma 2.1 and its proof,
\url{https://doi.org/10.1006/jcta.2000.3155}, use a Paley block and
an asymptotic-existence block in a short progression. The conclusion
below tracks the new sublogarithmic multiplier bound rather than
replacing it by an arbitrary fixed power. No novelty claim is made
pending a dedicated literature comparison.

\begin{theorem}[Conditional quantitative half-frame]
Assume GRH for primitive Dirichlet $L$-functions whose conductors are
powers of two, including conductor one. There are positive constants
$B,K$ such that every sufficiently large $n\equiv0\pmod4$ satisfies
\[
 \boxed{r(n)\ge\frac n2-
 B\sqrt n\exp\!\left(K\frac{\log n}{\log\log n}\right).}
\]
In particular the conditional deficit is $n^{1/2+o(1)}$ with a
specified subpower envelope, not merely $O_\epsilon(n^{1/2+\epsilon})$.
The theorem remains conditional even though its $\tau$ input is
unconditional.
\end{theorem}

\subsection{A uniform multiplier buffer removes circularity}
\begin{lemma}[Uniform power-of-two buffer]
There is an integer-valued function $T(n)\ge2$ such that
\[
 T(n)=O\!\left(\frac{\log n}{\log\log n}\right),
 \qquad \tau(u)\le T(n)\quad
 \text{for every positive odd }u\le n.
\]
For $Q(n)=2^{T(n)}$, every integer $1\le k\le n$ has a Hadamard
matrix of order $Q(n)k$.
\end{lemma}
\begin{proof}
Choose $A,u_0$ so $\tau(u)\le A\log u/\log\log u$ for odd
$u\ge u_0$, increasing $u_0$ beyond $e^e$. The function
$\log x/\log\log x$ is increasing there. The standard all-odd
multiplier construction ensures $\tau(u)<\infty$ also for the finitely
many smaller odd parts (for example, the binary-length complementary
pair construction supplies them). A fixed additive integer buffer
covers these smaller odd $u$, so
$T(n)=\lceil A\log n/\log\log n\rceil+B_0$ works for sufficiently
large $n$, with $B_0\ge2$. Write $k=2^vu$ with $u$ odd. Then
$\tau(u)\le T(n)$, and Sylvester doubling raises order
$2^{\tau(u)}u$ to $2^{T(n)+v}u=Q(n)k$.
\end{proof}

In particular $Q(n)=n^{o(1)}$. The buffer is chosen from $n$ before
the progression element or $k$ is selected. No bound on $\tau(k)$
depending on an as-yet-unknown construction is being assumed.

\subsection{The precise conditional number-theory input}
Put
\[
 \psi(x;Q,a)=\sum_{\substack{m\le x\\m\equiv a\pmod Q}}\Lambda(m).
\]
Under Dirichlet GRH, uniformly for reduced residues and, in particular,
for $Q\le\sqrt x$,
\[
 \psi(x;Q,a)=\frac{x}{\varphi(Q)}+
 O\bigl(\sqrt x(\log x)^2\bigr),
\]
with an absolute implied constant. A primary author source is
Montgomery--Vaughan, \emph{Multiplicative Number Theory I},
Theorem 13.7 and Corollary 13.8,
\url{https://personal.science.psu.edu/rcv4/personal/Publications/MNTI/17.0_pp_419_451_Conditional_estimates.pdf}.
Only characters induced by conductors dividing the power-of-two
modulus are used, so the stated restricted GRH hypothesis suffices.

For some absolute $D>0$, this implies the following elementary
short-interval consequence: whenever $x$ is sufficiently large,
$Q\ge4$ is a power of two, and
\[
 H=DQ\sqrt x(\log x)^2\le x/2,
\]
every reduced class modulo $Q$ contains a prime power in $(x-H,x]$.
Indeed subtract the two $\psi$ estimates. The main term is
$H/\varphi(Q)=2D\sqrt x(\log x)^2$, whereas both error terms
together are bounded by a fixed multiple of
$\sqrt x(\log x)^2$. Choose $D$ larger than that fixed multiple.
Positive $\psi$ difference supplies a prime power. An actual prime
is unnecessary for the Paley construction, so no prime-power
subtraction estimate is needed.

\subsection{Construction and deficit estimate}
\begin{proof}[Proof of the theorem]
Fix a large multiple of four $n$. Set $Q=2^{T(n)}$ and
\[
 x=n/2-1,\qquad a\equiv n-1\pmod Q,\qquad
 H=DQ\sqrt x(\log x)^2.
\]
The residue $a$ is odd, hence coprime to $Q$. Since $Q=n^{o(1)}$,
we eventually have $Q\le\sqrt{x-H}$ and $H\le x/2$.
The conditional interval estimate gives a prime power
\[
 x-H<\ell\le x,\qquad \ell\equiv n-1\pmod Q.
\]
Because $4\mid Q$ and $4\mid n$, necessarily $\ell\equiv3\pmod4$.
Paley's type-I construction therefore gives a Hadamard matrix of
order $b=\ell+1$. Also
\[
 k=\frac{n-\ell-1}{Q}\in\mathbb Z_{>0},\qquad k\le n.
\]
The uniform-buffer lemma gives a Hadamard matrix of order
$a_1=Qk=n-b$. Here $b\le n/2$ and $a_1\ge n/2$.
Take the first $b$ rows of each matrix and concatenate horizontally.
Their row Gram matrices add to
$a_1I_b+bI_b=nI_b$, producing a partial Hadamard matrix with
\[
 r(n)\ge b>n/2-H.
\]
Finally
\[
 H\ll Q\sqrt n(\log n)^2
   =\sqrt n\exp\!\left(O\!\left(\frac{\log n}{\log\log n}\right)\right),
\]
since $\log\log n=o(\log n/\log\log n)$. This proves the claim.
\end{proof}

The progression residue is $n-1\pmod Q$, which need not equal
$-1\pmod Q$: the target order need not be divisible by the buffer.
This distinction is essential when $Q$ exceeds the power of two
dividing $n$.

More quantitatively, if
$\tau(u)\le A\log u/\log\log u+O(1)$, the same proof gives,
for every $\eta>0$ and all sufficiently large target orders,
\[
 r(n)\ge\frac n2-
 \sqrt n\exp\!\left((A\log2+\eta)
                 \frac{\log n}{\log\log n}\right).
\]
All fixed constants and logarithmic factors are absorbed by the
$\eta$ margin. No numerical threshold or explicit value of $A$
is claimed here.

The verifier supplies unconditional finite examples of the exact
concatenation step: it searches for actual primes in the required
residue, constructs both Hadamard blocks at known orders, and checks
the full rectangular Gram matrix as integers. These calculations
neither establish GRH nor verify the eventual prime-distribution
estimate or the asymptotic constants.

\section*{Inputs from the preceding work}
The unconditional multiplier input is
$\tau(u)=O(\log u/\log\log u)$ for positive odd $u$.
Its proof and formalization are preserved in the main project.
Define $T(x)=\max(2,\max_{u\le x,\ u\text{ odd}}\tau(u))$.
Eventual monotonicity and a finite initial buffer give
$T(x)=O(\log x/\log\log x)$, and Sylvester doubling gives
a Hadamard matrix of order $2^{T(x)}v$ for every integer $1\le v\le x$.

\begin{theorem}[Spanning-slope input, proved in the eleventh pass]
\label{thm:design-spanning-slopes}
Let $h=2^s$, let $k$ be a known Hadamard order, and assume
$s\le h-k<h/2$. There is an exactly unbiased pairwise-independent
uniform law on $h^2$ atoms and $kh-1$ bits. It is a strict local
collision maximum in the full moment polytope, and has collision
$1/h^2$ and total variation distance exactly $1-1/k$ from the convex
hull of all order-$kh$ Hadamard laws. Its additive R\'enyi entropy
excess is $\log(h/k)$ for every entropy order.
\end{theorem}
The proof is in \path{../round11/global_repair.tex}, Sections 2--4:
finite-field affine slopes provide exact moments, and an omitted
additive-spanning slope set forces the support orthogonality graph
to consist of $h$ disjoint $h$-cliques. No small-repair assertion
or unproved Hadamard existence assumption is inserted here.
\section{Quantitative stability counterexamples from ordinary prime gaps}
\label{sec:stability-prime-gaps}
This section replaces the irrational-rotation subsequence with a construction
at every sufficiently large dyadic seed order. No prime progression is
needed. We use the established Baker--Harman--Pintz theorem that
$[x-x^{21/40},x]$ contains a prime for all sufficiently large $x$:
\emph{The difference between consecutive primes, II} (2001), Theorem 1,
\url{https://www.cs.umd.edu/~gasarch/BLOGPAPERS/BakerHarmanPintz.pdf}.
We do not claim $21/40$ is the best current prime-gap exponent.
The separate conditional improvement uses ordinary RH, not GRH for
arithmetic progressions. No new Hadamard orders are asserted.

\begin{lemma}[No residue-class restriction is needed]
\label{lem:stability-any-prime}
For every odd prime $p$, a Hadamard matrix of order $k=2(p+1)$ exists.
\end{lemma}
\begin{proof}
For $p\equiv1\pmod4$, this is classical Paley type II. For
$p\equiv3\pmod4$, Paley type I gives order $p+1$, which doubles to
$2(p+1)$. Both constructions originate in Paley,
\emph{On Orthogonal Matrices} (1933),
\url{https://doi.org/10.1002/sapm1933121311}.
\end{proof}

\begin{theorem}[An ordinary-prime-gap transfer]
\label{thm:stability-prime-transfer}
Suppose $0<\theta<1$ and every sufficiently large interval
$[x-x^\theta,x]$ contains a prime. For every sufficiently large power
of two $h$, there is an explicitly searchable known Hadamard order
$k<h$ with
\[
 h-k=\Theta(h^\theta),\qquad k/h\to1.
\]
At the known Hadamard order $n=kh$, there is an exact unbiased
pairwise-independent uniform $h^2$-atom law $\mu_h$ on $n-1$ bits,
a strict local collision maximum, for which
\begin{align*}
 C(\mu_h)&=1/h^2, & C_{\rm global}&=1/n,\\
 C(\mu_h)/C_{\rm global}&=1-\Theta(n^{-(1-\theta)/2}),&
 \operatorname{dist}_{\rm TV}(\mu_h,\mathcal H_n)
 &=1-1/k=1-\Theta(n^{-1/2}),\\
 1/n-C(\mu_h)&=\Theta(n^{-(3-\theta)/2}).
\end{align*}
All moments are exact; there is no approximate-isotropy assumption.
\end{theorem}
\begin{proof}
Let $X=h/2$, $A=\lceil X^\theta\rceil$, and $y=X-A$.
Choose the largest prime $p\le y$. The input theorem, applied at $y$,
gives $p\ge y-y^\theta$. Set $k=2(p+1)$, so
\[
 2A-2\le g:=h-k\le2A+2y^\theta-2\le4A-2.
\]
Thus $g=\Theta(h^\theta)$, and eventually
$\log_2h\le g<h/2$. The lemma supplies $H_k$, hence $H_n$ by tensoring
with Walsh order $h$.

Choose an omitted-slope set $R\subset\mathbb F_h\setminus\{0\}$ of
size $g$ containing an additive basis, and put
$P=\mathbb F_h\setminus R$. For uniform $A_0,B_0\in\mathbb F_h$, observe
all nonzero additive characters of each seed block
$x_t=A_0+tB_0$, $t\in P$, and the mixed bits $\chi_e(x_0+x_t)$,
$t\in P\setminus\{0\}$. This is the spanning-omitted-slopes construction
of Theorem~\ref{thm:design-spanning-slopes}: it has $h^2$ distinct atoms,
exact moments, and is a strict local maximum because two block character
spaces generate every seed character through products. Its orthogonality
graph is $h$ disjoint $h$-cliques; their order-$k$ block completions give
exact distance $1-1/k$ from all Hadamard mixtures. Collision is $1/h^2$,
whereas the known full optimum is $1/(kh)$.
Finally $n\asymp h^2$, and
\[
 1-\frac{C(\mu_h)}{1/n}=\frac gh,\qquad
 \frac1n-C(\mu_h)=\frac{g}{kh^2}.
\]
These identities give all rates.
\end{proof}

\begin{corollary}[Unconditional polynomial rates]
\label{cor:stability-bhp-rate}
Unconditionally the theorem holds with $\theta=21/40$. In particular
\[
 1-C(\mu_h)/(1/n)=\Theta(n^{-19/80}),\qquad
 1/n-C(\mu_h)=\Theta(n^{-99/80}),\qquad
 \operatorname{dist}_{\rm TV}(\mu_h,\mathcal H_n)=1-\Theta(n^{-1/2}).
\]
The order selection is a finite descending prime search below
$h/2-\lceil(h/2)^{21/40}\rceil$; the asymptotic theorem does not claim
a numerical value for the Baker--Harman--Pintz threshold.
\end{corollary}

\subsection{A separate improvement under ordinary RH}
Dudek proves on RH that for every $x\ge2$ there is a prime in
$(x-\frac4\pi\sqrt{x}\log x,x]$:
\emph{On the Riemann hypothesis and the difference between primes},
\url{https://arxiv.org/abs/1402.6417}. This is a primary, explicit
ordinary-prime result; no arithmetic-progression hypothesis is used.
\begin{corollary}[Conditional RH rates]\label{cor:stability-rh-rate}
Assuming RH, the same construction can be chosen at every sufficiently
large dyadic $h$ with
\[
 h-k=\Theta(\sqrt h\log h),\quad
 1-C(\mu_h)/(1/n)=\Theta(n^{-1/4}\log n),\quad
 1/n-C(\mu_h)=\Theta(n^{-5/4}\log n),
\]
while the exact distance remains $1-\Theta(n^{-1/2})$.
\end{corollary}
\begin{proof}
Use $A=\lceil\frac4\pi\sqrt X\log X\rceil$ and $y=X-A$ in the
preceding argument. The function $\frac4\pi\sqrt x\log x$ is increasing
for $x\ge2$, so the prime supplied at $y$ gives
$2A-2\le h-k\le4A-2$. The remaining proof is unchanged.
\end{proof}

\paragraph{Entropy and the scale of a possible stability theorem.}
Every Renyi entropy, including Shannon and min-entropy, of the uniform
source equals $\log h^2$. The minimum among feasible laws at known order
$n$ is $\log n$, by the atom bound and the Hadamard law. Thus all these
entropy excesses are simultaneously
$\log(h/k)=\Theta(n^{-19/80})$ in the unconditional construction,
and $\Theta(n^{-1/4}\log n)$ under RH, although distance tends to one.
These examples refute a dimension-free stability modulus based on relative
collision deficit or entropy excess. They also quantify absolute-deficit
counterexamples, at the exponents above. They do not exclude much smaller
dimension-dependent absolute collision windows, such as the earlier
order-$n^{-4}$ extraction condition. All positive existence inputs here
are classical known Hadamard constructions.

\begin{corollary}[Only sublinearly many extra atoms]
\label{cor:stability-support-excess}
The unconditional source support size satisfies
\[
 |\operatorname{supp}\mu_h|=n+\Theta(n^{61/80}).
\]
Under RH it can instead be $n+\Theta(n^{3/4}\log n)$, while its exact
distance from all Hadamard mixtures still tends to one.
\end{corollary}
\begin{proof}
The atom count is $h^2$ and $h^2-n=h(h-k)$. Substitute the corresponding
prime-gap scale and $n\asymp h^2$.
\end{proof}
\section{A transfer from existing Hadamard-order gap bounds}
\label{sec:stability-hadamard-gap-transfer}
The prime construction is an independent published-input route. More
generally, any uniform gap theorem for \emph{known Hadamard orders}
can be used directly, potentially giving substantially stronger rates.
The following statement deliberately isolates its hypothesis; applying a
particular gap bound requires proving or citing that bound separately.

\begin{theorem}[Uniform Hadamard-gap transfer]
\label{thm:stability-hadamard-gap-transfer}
Suppose $G$ is nonnegative and nondecreasing, $G(x)=o(x)$, and every
sufficiently large interval $[x-G(x),x]$ contains a known Hadamard order.
For dyadic $h=2^s$ put $A(h)=G(h)+s+1$. Then for every sufficiently large
$h$ there is a known Hadamard order $k$ with
\[
 2A(h)\le h-k\le3A(h).
\]
Consequently there is a uniform $h^2$-atom exactly isotropic strict local
collision maximum on $n-1$ bits, $n=kh$, with exact distance $1-1/k$
from all Hadamard mixtures and
\[
 \text{relative collision deficit}=\Theta(A(h)/h),\quad
 \text{absolute deficit}=\Theta(A(h)/h^3),\quad
 |\operatorname{supp}\mu|-n=\Theta(hA(h)).
\]
\end{theorem}
\begin{proof}
Apply the gap hypothesis at $x=h-2A(h)$. It gives
$k\in[x-G(x),x]$, with $G(x)\le G(h)\le A(h)$.
Thus $2A(h)\le h-k\le3A(h)$, and eventually
$s\le h-k<h/2$. Choose the omitted slopes to contain an additive basis
and apply Theorem~\ref{thm:design-spanning-slopes}. Its exact identities
are $1-C/(1/n)=(h-k)/h$, $1/n-C=(h-k)/(kh^2)$, and
$h^2-n=h(h-k)$. Since $k\sim h$, the stated estimates follow.
\end{proof}

For a polynomial gap input $G(x)=Cx^\theta$, $0<\theta<1$, the same
exponents as the prime-gap transfer result. For a subpolynomial input
\[
 G(x)=\exp\!\left(O\!\left(\frac{\log x}{\log\log x}\right)\right)
\]
with a nondecreasing majorant, the conclusion strengthens to
\begin{align*}
 \text{relative deficit}&=n^{-1/2+o(1)},&
 \text{absolute deficit}&=n^{-3/2+o(1)},\\
 |\operatorname{supp}\mu|&=n+n^{1/2+o(1)},&
 \operatorname{dist}_{\rm TV}(\mu,\mathcal H_n)&=1-\Theta(n^{-1/2}).
\end{align*}
These are conditional on the stated gap input, not on RH or any
arithmetic-progression conjecture. In particular, a gap theorem already
proved elsewhere in the project can be inserted here with its exact
scope and provenance. This transfer does not independently prove that
gap theorem or the Hadamard conjecture.
\section{Instantiating the stability transfer with the project's multiplier theorem}
\label{sec:stability-project-input}

The uniform Hadamard-gap input is already available in the current
project, not an additional prime-distribution assumption. Precisely,
\path{../../../papers/hadamard_multibase_final.tex}, the uniform-gap
theorem, proves that for each large integer $x$ there is
$d\le\exp(D\log x/\log\log x)$ such that every integer $y\le x$ has a
Hadamard order in $(y,y+d]$. Its formal counterpart is
\path{../../../lean/HadamardMultibase/CompleteResults.lean},
\path{Complete.quantitative_gap_bound}. We use the theorem as stated;
no files in that proof or its verification manifests are changed.

Here is an even more direct dyadic construction using just the
unconditional multiplier input and the uniform buffer lemma of the
arithmetic branch. For $h=2^s$ large choose
\[
 S(h)=\max\{T(h),\lceil\log_2s\rceil\},\qquad
 g=2^{S(h)},\qquad k=h-g.
\]
Since $T(h)=O(\log h/\log\log h)$, eventually
$S(h)\le s-2$, while $g\ge s$ and
\[
                 g=\exp(O(\log h/\log\log h))=h^{o(1)}.
\]
The integer $q=2^{s-S(h)}-1$ is positive odd and at most $h$.
Its admissible exponent is at most $T(h)\le S(h)$, so
\[
                    k=2^{S(h)}q
\]
is a known Hadamard order. This choice is uniform in all sufficiently
large dyadic $h$ and does not use GRH, prime gaps, or an
irrational-rotation subsequence.

\begin{corollary}[Unconditional near-square instability from the multiplier bound]
\label{cor:stability-project-rate}
At orders $n=kh$ obtained above, there are exactly isotropic uniform
laws on $m=h^2$ atoms, on $n-1$ bits, with
\[
 \begin{split}
 C(\mu)/(1/n)&=1-n^{-1/2+o(1)},\\
 1/n-C(\mu)&=n^{-3/2+o(1)},\\
 m&=n+n^{1/2+o(1)},\\
 \operatorname{dist}_{\rm TV}(\mu,\mathcal H_n)
       &=1-(1+o(1))n^{-1/2}.
 \end{split}
\]
They are strict local collision maxima. Their common additive
R\'enyi entropy excess is $n^{-1/2+o(1)}$.
Both $H_k$ and $H_n$ exist unconditionally from the project's
proved inputs.
\end{corollary}
\begin{proof}
The construction gives $s\le g<h/2$. Choose $g$ nonzero omitted
slopes containing an additive basis and apply the spanning-slope
theorem of the preceding pass. It yields exact isotropy, strict
local maximality, collision $1/h^2$, and exact distance $1-1/k$.
Tensoring $H_k$ and Walsh $H_h$ gives $H_n$.
Since $n\sim h^2$ and $g=h^{o(1)}$, the identities
\[
 1-\frac{C}{1/n}=\frac gh,\quad
 \frac1n-C=\frac{g}{kh^2},\quad
 m-n=hg,\quad
 H_\alpha(\mu)-\log n=\log(h/k)\sim g/h
\]
give all the assertions.
\end{proof}

Thus the project's own unconditional existence theorem gives a
stronger asymptotic stability counterexample than the independent
ordinary-prime-gap route, even assuming RH in that prime route.
The RH rate remains useful as an independently sourced comparison.
The positive all-order existence question is unchanged: these
counterexamples occur at orders where matrices are already known.
They concern full $n-1$-bit isotropy, not the punctured extraction input.
\section{Third completion moments with dependent triples included}
Fix an odd prime $p$, length $N=4p$, and a totally isotropic retained
space $G$ of codimension $h\ge2$ in a self-dual code. In the split
$2h$-dimensional quotient $G^\perp/G$, give each isotropic projective
point $\ell$ weight $w_\ell$ equal to the number of normalized sign
classes with nonzero quotient on that point. Let $b$ count the sign
classes in $G$. For a uniformly chosen completion code write
$Y=|S(C)|-b$. Put $D_j=p^{h-j}+1$.

For $t=1,2,3$ let $R_t$ be the sum of
$w_{\ell_1}w_{\ell_2}w_{\ell_3}$ over ordered triples whose span is
totally isotropic of dimension exactly $t$. Repetitions are allowed.
In particular $R_1=\sum_\ell w_\ell^3$. For $h=2$ necessarily
$R_3=0$; the definition of $D_3$ is then unnecessary.

\begin{theorem}[Exact third-moment certificate]
For $h\ge3$,
\[
 \mathbb E Y^3=\frac{R_1}{D_1}+
 \frac{R_2}{D_1D_2}+\frac{R_3}{D_1D_2D_3}.
\]
For $h=2$ omit the final term. If $\mathbb EY^2>0$, some completion
satisfies $|S(C)|\ge b+\lceil\mathbb EY^3/\mathbb EY^2\rceil$.
\end{theorem}
\begin{proof}
The number of generators in a split $2j$-space is
$L_j=2\prod_{i=1}^{j-1}(p^i+1)$ for $j\ge1$, and $L_0=1$.
A specified totally isotropic $t$-space lies in proportion
$L_{h-t}/L_h=\prod_{j=1}^tD_j^{-1}$ of the generators.
This includes $t=h$, where the final factor is $D_h=2$.
A triple not spanning an isotropic space cannot occur in a generator.
Expanding the cube of $Y$ and grouping by span rank proves the formula.
Finally $Y^3\le(\max Y)Y^2$ pointwise gives the certificate.
\end{proof}
Distinct points need not be linearly independent: three different
points on one isotropic projective line contribute to $R_2$, not
$R_3$. Omitting these triples would give an incorrect moment.

\subsection{An exact computation without cubic enumeration}
Let $\mathcal L$ be the totally isotropic projective lines. Then
\[
 R_2=\sum_{L\in\mathcal L}
 \left[\left(\sum_{\ell\in L}w_\ell\right)^3
              -\sum_{\ell\in L}w_\ell^3\right].
\]
Each rank-two triple spans a unique such line, so it is counted once.
Alternatively form the simple orthogonality graph on the weighted
points. Let $T_\triangle$ be the weighted sum over unordered triangles,
and $T_{\rm line}$ the weighted sum over unordered collinear triples.
For unordered orthogonal pairs $\{\ell,m\}$,
\[
 R_2=3\sum_{\{\ell,m\}}w_\ell w_m(w_\ell+w_m)+6T_{\rm line},
 \qquad R_3=6(T_\triangle-T_{\rm line}).
\]
Bitset intersections count graph triangles efficiently when all point
weights are one. For each pair, the other $p-1$ points on its line
are obtained by normalizing $x+ty$, $t\in\mathbb F_p^*$; dictionary
lookup counts collinear triples. Requiring the third index to exceed
both pair indices counts each once. The method uses quadratic pair
enumeration and avoids testing all ordered point triples.

\subsection{Known-order exact benchmarks}
For the first eight stored rows of the order-twenty terminal frame,
$b=9$, $h=2$, and
\[
 R_1=258,\quad R_2=7632,\quad R_3=0,\quad
 \mathbb EY^2=65,\quad\mathbb EY^3=679.
\]
Thus $|S(C)|\ge9+\lceil679/65\rceil=20$.
The second-moment bound gave only nineteen. Independent enumeration
of all twelve generators gives additional counts $3$ six times and
$11$ six times, directly checking the cubic moment.

For the first five stored rows, $b=5$, $h=5$, and the 1092 weighted
points all have weight one. The graph has 6497280 unordered triangles,
of which 1200 are collinear. Exact arithmetic gives
\[
 R_1=1092,\quad R_2=1118736,\quad R_3=38976480,
 \qquad\mathbb EY^3=\frac{995014}{28483}.
\]
As $\mathbb EY^2=14114/2191$, the ratio is
$497507/91741$, giving only $5+\lceil497507/91741\rceil=11$
total classes. This improves the second-moment lower bound of nine,
but does not reach the sufficient threshold seventeen.

\subsection{Scope of the moment hierarchy}
For any integer $r\ge2$ with positive denominator,
\[
 \max Y\ge\frac{\mathbb EY^r}{\mathbb EY^{r-1}}.
\]
These ratios are nondecreasing by Cauchy--Schwarz and tend to
$\max Y$ on the finite completion space. Since $Y$ is integer valued,
their ceilings eventually equal its exact maximum. This is a
classical moment principle, not a uniform bound on the number of
moments required or their computational cost. Supplying sufficiently
large finite-geometry counts remains the missing existence input in
unknown orders. No all-order positivity assertion follows.
\section{Rational dilations of arithmetic frames and an exact completion obstruction}
\label{sec:rational-dilations}

Concatenation produces the partial matrices used in the arithmetic
branch. It is useful to understand the remaining completion problem
as an explicit orthogonal dilation. The elementary calculation below
separates real or rational completion from the sign constraints.
No novelty is claimed for the linear algebra or the quaternion
sum-of-four-squares construction.

Let $a<b$ be positive multiples of four, let $A$ be an order-$a$
Hadamard matrix, and partition an order-$b$ Hadamard matrix as
$H_b=\binom{B}{E}$, with $B$ having $a$ rows. Put $d=b-a$ and $n=a+b$.
Then $P=[A\ B]$ is an $a$-row sign frame of width $n$.

\begin{theorem}[An exact residual parameterization]
\label{dil:parameterization}
The real matrices
\[
 Z=\begin{pmatrix}A&B\\B^T&D\end{pmatrix},\qquad ZZ^T=nI_n,
\]
are exactly those with
\[
 \boxed{D=\frac1b\left(-B^TA^TB+E^TVE\right),\qquad
                         VV^T=nI_d.}
\]
In particular a rational such $Z$ always exists. Its only possibly
nonsign entries are in the lower-right $b$-square block.
A sign-valued completion with this prescribed transpose corner
exists if and only if some such $V$ makes every entry of $D$ equal
to $+1$ or $-1$.
\end{theorem}
\begin{proof}
The top block already has Gram $nI_a$. The remaining equations are
\[
 AB+BD^T=0,\qquad B^TB+DD^T=nI_b.
\]
Conjugate $D$ by the orthogonal matrix
$W=\binom{B}{E}/\sqrt b$. The first equation forces
\[
 WDW^T=\begin{pmatrix}-A^T&F\\0&V\end{pmatrix}.
\]
The second requires its product with its transpose to be
$\operatorname{diag}(aI_a,nI_d)$.
Since $A^TA=aI_a$, this forces $FF^T=0$, hence $F=0$,
and leaves exactly $VV^T=nI_d$. Transforming back gives the formula.

For rational existence write $n=u_0^2+u_1^2+u_2^2+u_3^2$ by
Lagrange's four-square theorem. The quaternion multiplication matrix
\[
 Q=\begin{pmatrix}
 u_0&u_1&u_2&u_3\\
 -u_1&u_0&-u_3&u_2\\
 -u_2&u_3&u_0&-u_1\\
 -u_3&-u_2&u_1&u_0
 \end{pmatrix}
\]
satisfies $QQ^T=nI_4$. Since $d$ is a multiple of four, take $V$
block diagonal with $d/4$ copies of $Q$. It is integral, and the
displayed formula gives a rational $D$. The sign-valued criterion
is then simply the definition of a Hadamard matrix with these
prescribed three blocks.
\end{proof}

The free real parameter lies in a scaled $O(d)$, rather than $O(n)$.
This is a reduction of the completion variables, not a proof of
flatness. In particular rational scaled orthogonality supplies no
general sign-rounding theorem.

\subsection{A structured order-twenty obstruction and its smallest repair}
Take $A=H_8$ in the standard Sylvester order, and let $B$ be the
first eight rows of the normalized Paley type-II $H_{12}$ constructed
from the prime five. The exact frame $[A\ B]$ has eight orthogonal
rows of length twenty. It cannot be completed to a Hadamard matrix
while preserving all eight rows.

Here is a finite exhaustive certificate of that assertion.
Its row span modulo five has dimension eight. The four-dimensional
hyperbolic quotient therefore has exactly twelve self-dual
extensions, the two rulings of the preceding pass's $6$-square grid.
The checker reconstructs all twelve codes and exhausts each
systematic sign assignment. Each code has exactly twelve sign
classes, and those classes are real orthogonal because the
prescribed eight rows fix their parity. Thus the exact maximum
orthogonal frame preserving these eight rows is twelve, not twenty.
The row residues of any full Hadamard completion would lie in one
of these self-dual extensions, contradicting its sign count.

Deleting one prescribed row does permit a full completion.
The checker enumerates modular candidates orthogonal to the retained
seven rows, makes an explicit self-dual neighbor pivot, and obtains
twenty pairwise orthogonal sign classes preserving exactly seven of
the eight originals. The certificate and row index are in
\path{dilations/checks.json}. The minimum original-row loss for a
full completion is consequently exactly one.
This is a known-order example; it proves neither a uniform bounded
repair theorem nor failure of the arithmetic construction at other orders.

The continuous family is nonempty despite the sign obstruction.
Using the quaternion parameters $(4,2,0,0)$ gives $VV^T=20I_4$.
The checker verifies the full rational $20$-square Gram identity,
with lower-right entries ranging from $-2$ to $8/3$ and
\[
                    \sum_{ij}D_{ij}^4-144=\frac{6976}{27}.
\]
Rounding its entries to signs does not produce a Hadamard matrix.

\subsection{A quantitative flatness gap when the prescribed rows cannot complete}
\begin{proposition}[A rounding threshold for real completions]
\label{dil:flatness-gap}
Let $n$ be a positive multiple of four and let $Z$ be any real matrix
with $ZZ^T=nI_n$. Define
\[
 \Delta=\sum_{ij}(Z_{ij}^2-1)^2=\sum_{ij}Z_{ij}^4-n^2.
\]
If
\[
 \Delta<(\sqrt{n+2}-\sqrt n)^2,
\]
rounding each entry to a sign gives a Hadamard matrix.
Consequently if some prescribed sign rows cannot be completed,
every real completion preserving them has
$\Delta\ge(\sqrt{n+2}-\sqrt n)^2>1/(n+2)$.
\end{proposition}
\begin{proof}
Let $H$ be the entrywise sign rounding (choose $+1$ at zero)
and write $F=H-Z$. The scalar inequality
$(|z|-1)^2\le(z^2-1)^2$ gives $\|F\|_F\le\sqrt\Delta$.
Since $\|Z\|_{\rm op}=\sqrt n$,
\[
 \|HH^T-nI_n\|_{\rm op}
 \le2\sqrt n\,\sqrt\Delta+\Delta<2.
\]
The diagonal is zero, and every off-diagonal entry of this error
is an even integer, since $n$ is even and $H$ is a sign matrix.
The operator norm bounds every entry in magnitude, so all entries
vanish. Prescribed sign rows remain unchanged by rounding.
Finally
$(\sqrt{n+2}-\sqrt n)^2
=4/(\sqrt{n+2}+\sqrt n)^2>1/(n+2)$.
\end{proof}

For the order-twenty frame this gives an explicit positive quartic
gap for every continuous completion preserving the eight rows.
A continuous optimizer cannot remove that gap without changing
prescribed rows. This pinpoints a real limitation of feeding an
arithmetic partial frame into a monotone completion method.

\section{What remains unresolved}
The arithmetic construction reaches roughly half the required rows,
and the rational dilation solves scaled orthogonality without solving
its sign constraints. The structured finite example requires changing
a prescribed row before reaching a full matrix.
The third-moment method improves one weak count certificate only to
eleven, still below its completion threshold. No uniform high-moment
or count-surplus theorem is proved.

The stronger instability family shows that a fixed relative
collision or entropy tolerance cannot guarantee a small repair across
dimensions. It does not exclude the dimension-dependent punctured
extraction route. A full-conjecture proof still needs an actual
all-order construction or a global exchange theorem that reaches its
target without assuming the desired matrix exists.
The checks validate finite arithmetic certificates, not GRH or the
asymptotic constants. No Lean code, commits or submissions were made.
\end{document}
